% ============================================================
\chapter{GNN Variants (GAT, GIN, GraphSAGE)}
\label{ch:variants}
% ============================================================

The GCN of Kipf and Welling (2017) opened the way, but its limitations quickly became apparent: isotropic aggregation, expressivity capped at the 1-Weisfeiler-Leman test, and inability to handle large graphs inductively. In response, a wave of alternative architectures emerged: \emph{Graph Attention Networks} (GAT) introduce attention mechanisms to weight neighbours; \emph{Graph Isomorphism Networks} (GIN) maximise theoretical expressivity; \emph{GraphSAGE} enables inductive learning through neighbourhood sampling. This chapter explores these variants and the trade-offs they embody.

\section{Beyond GCN: Motivations}

The GCN of Kipf and Welling has several limitations:
\begin{itemize}
    \item \textbf{Isotropic} aggregation: all neighbors are treated equally
          (weighted only by degree).
    \item \textbf{Limited expressivity}: strict equivalence to 1-WL.
    \item \textbf{Scalability}: requires the full adjacency matrix.
\end{itemize}

This chapter presents three major variants addressing these limitations.

\section{Graph Attention Networks (GAT)}

\subsection{Attention on Graphs}

\begin{definition}[GAT layer --- Veli\v{c}kovi\'c et al., 2018]
The GAT layer uses an \textbf{attention} mechanism to weight neighbor
contributions:
\[
    \mathbf{h}_v^{(\ell+1)} = \sigma\!\left(\sum_{u \in \mathcal{N}(v) \cup \{v\}}
    \alpha_{vu}^{(\ell)}\, \mathbf{W}^{(\ell)} \mathbf{h}_u^{(\ell)}\right)
\]
where the attention coefficients $\alpha_{vu}$ are computed by:
\begin{align}
    e_{vu}^{(\ell)} &= \text{LeakyReLU}\!\left(\mathbf{a}^{\top}\!
    \left[\mathbf{W}\mathbf{h}_v^{(\ell)} \| \mathbf{W}\mathbf{h}_u^{(\ell)}\right]\right) \\
    \alpha_{vu}^{(\ell)} &= \frac{\exp(e_{vu}^{(\ell)})}{\sum_{w \in \mathcal{N}(v) \cup \{v\}}
    \exp(e_{vw}^{(\ell)})}
\end{align}
where $\mathbf{a} \in \R^{2d'}$ is the learned attention vector and $\|$
denotes concatenation.
\end{definition}

\begin{remark}[Multi-head attention]
To stabilize training, GAT uses $K$ independent attention heads:
\[
    \mathbf{h}_v^{(\ell+1)} = \Big\|_{k=1}^{K} \sigma\!\left(\sum_{u \in \mathcal{N}(v)}
    \alpha_{vu}^{k} \mathbf{W}^k \mathbf{h}_u^{(\ell)}\right)
\]
The last layer typically uses averaging instead of concatenation.
\end{remark}

\begin{intuition}[title=\textbf{Why attention?}]
Attention allows the model to learn \textbf{which neighbors matter} for each
node, unlike GCN where weighting is fixed by graph structure. This is
especially useful when:
\begin{itemize}
    \item Neighbors have \textbf{varying relevance}.
    \item The graph contains \textbf{noise} (uninformative edges).
    \item Relationships are \textbf{asymmetric} (directed graphs).
\end{itemize}
\end{intuition}

\begin{codeblock}[title=\textbf{GAT implementation with PyTorch Geometric}]
\begin{minted}{python}
import torch
import torch.nn.functional as F
from torch_geometric.nn import GATConv
from torch_geometric.datasets import Planetoid

class GAT(torch.nn.Module):
    def __init__(self, in_channels, hidden_channels, out_channels,
                 heads=8, dropout=0.6):
        super().__init__()
        self.conv1 = GATConv(in_channels, hidden_channels,
                             heads=heads, dropout=dropout)
        self.conv2 = GATConv(hidden_channels * heads, out_channels,
                             heads=1, concat=False, dropout=dropout)
        self.dropout = dropout

    def forward(self, x, edge_index):
        x = F.dropout(x, p=self.dropout, training=self.training)
        x = F.elu(self.conv1(x, edge_index))
        x = F.dropout(x, p=self.dropout, training=self.training)
        x = self.conv2(x, edge_index)
        return x

dataset = Planetoid(root='/tmp/Cora', name='Cora')
data = dataset[0]
model = GAT(dataset.num_features, 8, dataset.num_classes, heads=8)
optimizer = torch.optim.Adam(model.parameters(), lr=0.005, weight_decay=5e-4)

for epoch in range(200):
    model.train()
    optimizer.zero_grad()
    out = model(data.x, data.edge_index)
    loss = F.cross_entropy(out[data.train_mask], data.y[data.train_mask])
    loss.backward()
    optimizer.step()

model.eval()
pred = model(data.x, data.edge_index).argmax(dim=1)
acc = (pred[data.test_mask] == data.y[data.test_mask]).float().mean()
print(f"GAT test accuracy: {acc:.4f}")
\end{minted}
\end{codeblock}

\subsection{GATv2: Dynamic Attention}

\begin{definition}[GATv2 --- Brody et al., 2022]
GATv2 fixes a limitation of GAT where attention is \textbf{static} (the
ranking of neighbors does not depend on the query node). GATv2 uses:
\[
    e_{vu} = \mathbf{a}^\top \text{LeakyReLU}\!\left(\mathbf{W}\!\left[
    \mathbf{h}_v \| \mathbf{h}_u\right]\right)
\]
The nonlinearity is applied \textbf{before} the dot product with $\mathbf{a}$,
making the attention truly \textbf{dynamic}.
\end{definition}

\section{Graph Isomorphism Network (GIN)}

\subsection{Maximizing Expressive Power}

\begin{theorem}[GIN --- Xu et al., 2019]
For an MPNN to achieve the maximum power of the 1-WL test, the following
conditions are necessary and sufficient:
\begin{enumerate}
    \item The aggregation must be \textbf{injective} on multisets.
    \item The update function must be \textbf{injective}.
\end{enumerate}
\end{theorem}

\begin{definition}[GIN layer]
The GIN layer computes:
\[
    \mathbf{h}_v^{(\ell+1)} = \text{MLP}^{(\ell)}\!\left((1 + \epsilon^{(\ell)})
    \cdot \mathbf{h}_v^{(\ell)} + \sum_{u \in \mathcal{N}(v)}
    \mathbf{h}_u^{(\ell)}\right)
\]
where $\epsilon$ is a learned or fixed parameter.
\end{definition}

\begin{proposition}[Injectivity of sum]
Among classical aggregations (sum, mean, max), only \textbf{sum} is injective
on multisets:
\begin{itemize}
    \item $\text{max}\{1, 1, 2\} = \text{max}\{1, 2\} = 2$ (information loss).
    \item $\text{mean}\{1, 1\} = \text{mean}\{1\} = 1$ (cardinality loss).
    \item $\text{sum}\{1, 1, 2\} = 4 \neq \text{sum}\{1, 2\} = 3$ (injective).
\end{itemize}
\end{proposition}

\begin{codeblock}[title=\textbf{GIN implementation}]
\begin{minted}{python}
import torch
import torch.nn as nn
import torch.nn.functional as F
from torch_geometric.nn import GINConv, global_add_pool
from torch_geometric.datasets import TUDataset
from torch_geometric.loader import DataLoader

class GIN(nn.Module):
    def __init__(self, num_features, hidden_dim, num_classes, num_layers=5):
        super().__init__()
        self.convs = nn.ModuleList()
        self.bns = nn.ModuleList()
        for i in range(num_layers):
            in_dim = num_features if i == 0 else hidden_dim
            mlp = nn.Sequential(
                nn.Linear(in_dim, hidden_dim),
                nn.ReLU(),
                nn.Linear(hidden_dim, hidden_dim),
            )
            self.convs.append(GINConv(mlp, train_eps=True))
            self.bns.append(nn.BatchNorm1d(hidden_dim))
        self.classifier = nn.Linear(hidden_dim, num_classes)

    def forward(self, data):
        x, edge_index, batch = data.x, data.edge_index, data.batch
        for conv, bn in zip(self.convs, self.bns):
            x = F.relu(bn(conv(x, edge_index)))
        x = global_add_pool(x, batch)
        return self.classifier(x)

dataset = TUDataset(root='/tmp/MUTAG', name='MUTAG')
loader = DataLoader(dataset, batch_size=32, shuffle=True)
model = GIN(dataset.num_features, 64, dataset.num_classes)
print(f"Parameters: {sum(p.numel() for p in model.parameters()):,}")
\end{minted}
\end{codeblock}

\section{GraphSAGE: Inductive Learning}

\subsection{Motivation: Scalability}

\begin{definition}[GraphSAGE --- Hamilton et al., 2017]
GraphSAGE (\textit{SAmple and aggreGatE}) uses \textbf{neighborhood sampling}
for scalability:
\[
    \mathbf{h}_v^{(\ell+1)} = \sigma\!\left(\mathbf{W}^{(\ell)} \cdot
    \text{CONCAT}\!\left(\mathbf{h}_v^{(\ell)},\;
    \text{AGG}(\{\mathbf{h}_u^{(\ell)} : u \in \mathcal{S}(v)\})\right)\right)
\]
where $\mathcal{S}(v) \subseteq \mathcal{N}(v)$ is a sampled subset of
neighbors (typically $\abs{\mathcal{S}(v)} \leq k$).
\end{definition}

\begin{proposition}[Advantages of GraphSAGE]
\begin{enumerate}
    \item \textbf{Inductive}: can generalize to unseen nodes.
    \item \textbf{Scalable}: complexity controlled by sample count.
    \item \textbf{Mini-batch}: compatible with mini-batch training.
\end{enumerate}
\end{proposition}

\begin{algorithme}[title=\textbf{Subgraph sampling for GraphSAGE}]
\begin{enumerate}
    \item For each target node $v$, sample $k_1$ direct neighbors.
    \item For each of those neighbors, sample $k_2$ neighbors (2-hop).
    \item Form the induced subgraph and propagate messages.
    \item Per-batch complexity is $\mathcal{O}(B \cdot k_1 \cdot k_2)$
          instead of $\mathcal{O}(n)$.
\end{enumerate}
\end{algorithme}

\begin{codeblock}[title=\textbf{GraphSAGE with neighborhood sampling}]
\begin{minted}{python}
import torch
import torch.nn.functional as F
from torch_geometric.nn import SAGEConv
from torch_geometric.loader import NeighborLoader
from torch_geometric.datasets import Planetoid

dataset = Planetoid(root='/tmp/Cora', name='Cora')
data = dataset[0]

class GraphSAGE(torch.nn.Module):
    def __init__(self, in_channels, hidden_channels, out_channels):
        super().__init__()
        self.conv1 = SAGEConv(in_channels, hidden_channels)
        self.conv2 = SAGEConv(hidden_channels, out_channels)

    def forward(self, x, edge_index):
        x = F.relu(self.conv1(x, edge_index))
        x = F.dropout(x, p=0.5, training=self.training)
        x = self.conv2(x, edge_index)
        return x

train_loader = NeighborLoader(
    data,
    num_neighbors=[25, 10],
    batch_size=256,
    input_nodes=data.train_mask,
)

model = GraphSAGE(dataset.num_features, 64, dataset.num_classes)
optimizer = torch.optim.Adam(model.parameters(), lr=0.01)

for epoch in range(20):
    model.train()
    for batch in train_loader:
        optimizer.zero_grad()
        out = model(batch.x, batch.edge_index)[:batch.batch_size]
        loss = F.cross_entropy(out, batch.y[:batch.batch_size])
        loss.backward()
        optimizer.step()
\end{minted}
\end{codeblock}

\section{Variant Comparison}

\begin{center}
\begin{tabular}{lcccc}
\toprule
\textbf{Property} & \textbf{GCN} & \textbf{GAT} & \textbf{GIN} & \textbf{SAGE} \\
\midrule
Attention & No & Yes & No & No \\
1-WL expressivity & $<$ & $<$ & $=$ & $<$ \\
Inductive & No & Yes & Yes & Yes \\
Scalable (mini-batch) & Hard & Moderate & Moderate & Yes \\
Edge features & No & Yes (v2) & No & No \\
\bottomrule
\end{tabular}
\end{center}

\section{Advanced Architectures}

\subsection{PNA --- Principal Neighbourhood Aggregation}

\begin{definition}[PNA --- Corso et al., 2020]
PNA combines \textbf{multiple aggregators} and \textbf{scalers}:
\[
    \mathbf{h}_v^{(\ell+1)} = U\!\left(\mathbf{h}_v^{(\ell)},
    \underset{\substack{\oplus \in \{\text{sum, mean,}\\\text{max, std}\}}}
    {\Big\|}
    \underset{s \in \{1, \sqrt{d}, 1/\sqrt{d}\}}{\Big\|}
    s \cdot \bigoplus_{u \in \mathcal{N}(v)} M(\mathbf{h}_v^{(\ell)},
    \mathbf{h}_u^{(\ell)})\right)
\]
The concatenation of $4 \times 3 = 12$ channels captures different aspects
of the neighborhood distribution.
\end{definition}

\subsection{Residual Connections for Deep GNNs}

\begin{definition}[GCNII --- Chen et al., 2020]
GCNII adds an \textbf{initial residual connection} and \textbf{identity
mapping} to avoid over-smoothing:
\[
    \mathbf{H}^{(\ell+1)} = \sigma\!\left(\left((1-\alpha_\ell)\hat{\mathbf{P}}
    \mathbf{H}^{(\ell)} + \alpha_\ell \mathbf{H}^{(0)}\right)
    \left((1-\beta_\ell)\mathbf{I} + \beta_\ell \mathbf{W}^{(\ell)}\right)\right)
\]
where $\hat{\mathbf{P}} = \hat{\mathbf{D}}^{-1/2}\hat{\mathbf{A}}\hat{\mathbf{D}}^{-1/2}$
and $\alpha_\ell, \beta_\ell$ are hyperparameters.
\end{definition}

\section*{Exercises}

\begin{exercise}[Experimental comparison]
Compare GCN, GAT, GIN, and GraphSAGE on the Cora dataset:
\begin{enumerate}
    \item Implement all four models with the same parameter count.
    \item Train for 200 epochs with the same hyperparameters.
    \item Report test accuracy and loss curves.
    \item Analyze performance differences.
\end{enumerate}
\end{exercise}

\begin{exercise}[GAT attention visualization]
\begin{enumerate}
    \item Train a GAT on Cora and extract attention coefficients $\alpha_{vu}$.
    \item Visualize attention weights on a subgraph of 20 nodes.
    \item Are attention coefficients correlated with node labels?
\end{enumerate}
\end{exercise}

\begin{exercise}[GIN expressive power]
\begin{enumerate}
    \item Construct two graphs that GCN cannot distinguish but GIN can.
    \item Verify experimentally.
    \item Construct two graphs that even GIN cannot distinguish (1-WL failure).
\end{enumerate}
\end{exercise}

\begin{exercise}[Mini-batch training]
Implement mini-batch training with neighborhood sampling for a GCN (without
using \texttt{NeighborLoader}).
\begin{enumerate}
    \item Implement the sampling strategy.
    \item Compare training speed with full-batch mode.
    \item Evaluate the impact of the number of sampled neighbors on accuracy.
\end{enumerate}
\end{exercise}
